\documentclass[11pt,a4paper]{article}

% --- 极简中文支持（pdfLaTeX 直接兼容，无需切换编译器）---
\usepackage{CJKutf8} % pdfLaTeX 下最稳定的中文支持，无需 XeLaTeX

% --- 基础宏包（清理冗余，确保无冲突）---
\usepackage{amsmath, amsfonts, amssymb} % 数学公式（移除重复的 amsmath）
\usepackage{graphicx} % 插图（已注释调用）
\usepackage{geometry} % 页面布局
\usepackage{listings} % 代码块
\usepackage{xcolor} % 颜色
\usepackage{hyperref} % 超链接
\usepackage{booktabs} % 表格
\usepackage{caption}
\usepackage{float}

% 页面设置
\geometry{left=2.5cm, right=2.5cm, top=3cm, bottom=3cm}

% 代码风格设置
\lstset{
    basicstyle=\ttfamily\small,
    keywordstyle=\color{blue},
    commentstyle=\color{green!50!black},
    stringstyle=\color{red},
    frame=single,
    breaklines=true,
    captionpos=b
}

% --- 标题部分 ---
\title{计算机图形学实验报告：实验三 Poisson Image Editing}
\author{姓名：宁尚哲\quad 学号：PB24000099}
\date{\today}

\begin{document}
% 开启中文环境（整个文档生效）
\begin{CJK*}{UTF8}{gbsn}

\maketitle

\begin{abstract}
    本实验基于 Framework2D 框架，实现了经典的 Poisson Image Editing 算法。通过将图像克隆问题转化为离散泊松方程的求解，实现了无缝的图像融合（Seamless Cloning）。实验中利用 Eigen 库构建并求解大规模稀疏线性方程组，并针对交互体验实现了矩阵预分解的实时更新策略。此外，实验扩展支持了 Mixing Gradients 模式，以及基于扫描线算法的复杂不规则区域选择。实验结果证明，该方法能有效消除图像克隆时的边界断裂感，使融合结果在视觉上更加自然。
\end{abstract}

\section{引言}
在图像编辑中，传统的“剪切-粘贴”操作（Paste）往往会在目标图像中留下明显的边界接缝，这是由于源图像与目标图像背景颜色不一致导致的。Poisson Image Editing 提出在梯度域进行编辑：保持源区域内部的梯度信息（细节），同时通过求解偏微分方程使区域边界平滑过渡到目标背景颜色。

\section{算法原理}
\subsection{泊松方程与变分模型}
图像融合问题可建模为在已知边界条件下的泛函极小化问题。对于目标区域 $\Omega$，我们希望寻找一个函数 $f$，使得其梯度 $\nabla f$ 尽可能接近引导矢量场 $\mathbf{v}$（通常取自源图像 $\nabla g$）：
\begin{equation}
    \min_f \iint_{\Omega} |\nabla f - \mathbf{v}|^2 \, dA \quad \text{with} \quad f|_{\partial \Omega} = f^*|_{\partial \Omega}
\end{equation}
其中 $f^*$ 是目标图像，$\mathbf{v}$ 是引导矢量场，$\partial \Omega$ 是区域边界。这个能量泛函的含义是：在满足边界条件的前提下，使融合图像的梯度场尽可能接近引导场。

\subsubsection{变分原理与Euler-Lagrange方程}
根据变分法原理，上述泛函极小化问题对应的 Euler-Lagrange 方程为：
\begin{equation}
    \Delta f = \text{div} \mathbf{v} \quad \text{over} \quad \Omega, \quad \text{with} \quad f|_{\partial \Omega} = f^*|_{\partial \Omega}
\end{equation}
其中 $\Delta f = \frac{\partial^2 f}{\partial x^2} + \frac{\partial^2 f}{\partial y^2}$ 是拉普拉斯算子，$\text{div} \mathbf{v} = \frac{\partial v_x}{\partial x} + \frac{\partial v_y}{\partial y}$ 是散度算子。

\subsubsection{引导矢量场的选择}
\begin{itemize}
    \item \textbf{Seamless Cloning}：取 $\mathbf{v} = \nabla g$，即使用源图像的梯度场，保持源图像的细节纹理。
    \item \textbf{Mixing Gradients}：取 $\mathbf{v}(p) = \begin{cases} \nabla f^*(p), & \text{if } |\nabla f^*(p)| > |\nabla g(p)| \\ \nabla g(p), & \text{otherwise} \end{cases}$，即在每个像素点选择梯度绝对值较大的那个。
\end{itemize}

\subsection{离散化求解}
在离散网格上，对于每个像素 $p \in \Omega$，泊松方程可写为：
\begin{equation}
    |N_p| f_p - \sum_{q \in N_p \cap \Omega} f_q = \sum_{q \in N_p} v_{pq} + \sum_{q \in N_p \cap \partial\Omega} f^*_q
\end{equation}
其中 $N_p$ 为像素 $p$ 的 4-邻域，$v_{pq}$ 是引导矢量场在方向 $\overrightarrow{pq}$ 上的分量。

\subsubsection{有限差分格式}
使用中心差分近似拉普拉斯算子：
\begin{equation}
    \Delta f(x,y) \approx \frac{f(x+1,y) + f(x-1,y) + f(x,y+1) + f(x,y-1) - 4f(x,y)}{h^2}
\end{equation}
其中 $h$ 是网格间距（取 $h=1$）。对于内部点，离散化后的泊松方程为：
\begin{equation}
    4f(x,y) - f(x+1,y) - f(x-1,y) - f(x,y+1) - f(x,y-1) = \text{div} \mathbf{v}(x,y)
\end{equation}

\subsubsection{边界条件处理}
对于边界点，其值固定为目标图像的对应像素值。在离散化时，将边界项移到方程右侧：
\begin{equation}
    4f(x,y) - \sum_{q \in N_p \cap \Omega} f_q = \sum_{q \in N_p \cap \partial\Omega} f^*_q + \sum_{q \in N_p} v_{pq}
\end{equation}

\subsubsection{稀疏线性方程组}
将所有内部点的方程联立，得到稀疏线性方程组 $A\mathbf{x} = \mathbf{b}$，其中：
\begin{itemize}
    \item $A$ 是 $n \times n$ 的稀疏矩阵（$n$ 是内部点数量），每行最多5个非零元素。
    \item $\mathbf{x}$ 是待求解的像素值向量。
    \item $\mathbf{b}$ 是右侧向量，包含引导场散度和边界条件。
\end{itemize}

\section{算法实现细节}
\subsection{复杂区域掩码生成}
实验实现了基于扫描线（Scan-line）的算法来处理不规则多边形和自由绘制区域。

\subsubsection{扫描线算法原理}
扫描线算法是多边形光栅化的经典方法，其核心思想是：
\begin{enumerate}
    \item 构建边表（Edge Table）：将多边形的每条边按照其最小y坐标分类存储。
    \item 活动边表（Active Edge Table）：维护当前扫描线与多边形相交的所有边。
    \item 扫描线填充：从最小y坐标到最大y坐标逐行扫描，利用交点对填充内部像素。
\end{enumerate}

\subsubsection{数据结构设计}
在 \texttt{polygon.cpp} 中，定义了边的数据结构：
\begin{itemize}
    \item \texttt{Edge}：包含边的最大y坐标、当前x坐标和斜率倒数。
    \item \texttt{edge\_table}：映射表，键为y坐标，值为该y坐标处开始的边列表。
    \item \texttt{aet}（活动边表）：当前扫描线上的所有边。
\end{itemize}

\subsubsection{算法流程}
\begin{enumerate}
    \item 遍历多边形顶点，构建边表，计算每条边的斜率倒数 $\frac{1}{k} = \frac{\Delta x}{\Delta y}$。
    \item 对于每条扫描线 $y$，将边表中 $y$ 处的边加入活动边表。
    \item 移除活动边表中 $y_{\max} \leq y$ 的边。
    \item 对活动边表按x坐标排序，利用相邻交点对填充像素。
    \item 更新每条边的x坐标：$x_{\text{new}} = x_{\text{old}} + \frac{1}{k}$。
\end{enumerate}

\subsubsection{掩码生成}
在 \texttt{source\_image\_widget.cpp} 的 \texttt{update\_selected\_region()} 函数中：
\begin{itemize}
    \item 调用形状对象的 \texttt{get\_interior\_pixels()} 方法获取内部像素坐标。
    \item 清空掩码图像，将选中像素标记为255。
    \item 生成的掩码用于后续的泊松方程求解。
\end{itemize}

\subsection{方程组构建与加速策略}
实验的核心算法逻辑位于 \texttt{seamless\_clone.cpp}。

\subsubsection{坐标映射机制}
对于不规则区域，采用索引映射机制将二维坐标映射到一维索引：
\begin{itemize}
    \item \texttt{selected\_pixels\_}：存储所有选中像素的坐标 $(x, y)$，按顺序编号 $0, 1, 2, \ldots, n-1$。
    \item \texttt{coord\_to\_idx\_}：二维坐标到一维索引的映射，\texttt{coord\_to\_idx\_[y][x]} 返回像素 $(x,y)$ 的索引，未选中区域为 -1。
    \item 边界框优化：只遍历边界框内的像素，减少不必要的计算。
\end{itemize}

\subsubsection{系数矩阵构建}
在 \texttt{fill\_coefficient()} 函数中，为每个像素构建方程：
\begin{itemize}
    \item \textbf{内部点}：对角线元素为4，邻居元素为-1，右侧为引导场散度。
    \item \textbf{边界点}：对角线元素为1，其他为0，右侧为目标图像像素值。
    \item 使用 \texttt{Eigen::Triplet} 存储非零元素，构建稀疏矩阵。
\end{itemize}

\subsubsection{右侧向量计算}
\begin{itemize}
    \item \textbf{Seamless Cloning}：$B_i = 4g(x,y) - g(x+1,y) - g(x-1,y) - g(x,y+1) - g(x,y-1)$。
    \item \textbf{Mixing Gradients}：比较源图像和目标图像的梯度，选择绝对值较大的作为引导场。
    \item 对于边界邻居，将目标图像像素值加到右侧向量中。
\end{itemize}

\subsubsection{稀疏矩阵预分解}
\begin{enumerate}
    \item \textbf{预分解时机}：在 \texttt{target\_image\_widget.cpp} 的 \texttt{mouse\_click\_event()} 中，首次点击时调用 \texttt{precompute()}。
    \item \textbf{分解方法}：使用 \texttt{Eigen::SimplicialLDLT}，适用于对称正定矩阵，Poisson方程的系数矩阵满足此条件。
    \item \textbf{性能提升}：预分解后，每次拖动只需重新计算右侧向量 $B$ 并执行回代，速度提升约10-50倍。
\end{enumerate}

\subsubsection{实时编辑实现}
\begin{itemize}
    \item \textbf{快速求解}：在 \texttt{solve\_fast()} 中，使用预分解的求解器快速求解。
    \item \textbf{偏移更新}：在 \texttt{mouse\_move\_event()} 中，调用 \texttt{update\_offset()} 更新偏移量。
    \item \textbf{结果应用}：将解裁剪到 $[0, 255]$ 范围，设置到目标图像中。
\end{itemize}

\subsubsection{求解器选择}
实验中对比了多种稀疏求解器：
\begin{itemize}
    \item \textbf{SimplicialLDLT}：适用于对称正定矩阵，速度快，稳定性好。
    \item \textbf{SparseLU}：适用于一般稀疏矩阵，但速度较慢。
    \item \textbf{ConjugateGradient}：迭代法，适用于大规模矩阵，但收敛性依赖于初始值。
\end{itemize}
最终选择 \texttt{SimplicialLDLT}，因为Poisson方程的系数矩阵是对称正定的，且预分解后求解速度最快。

\section{实验结果与分析}

\subsection{模式对比实验}
本节对比了三种模式在同一场景下的效果：
\begin{figure}[H]
    \centering
    % 注释图片，避免路径错误
     \includegraphics[width=\textwidth]{figs/paste_result.png}
     \caption{直接粘贴：边界处颜色突变，产生明显的接缝}
\end{figure}

\begin{figure}[H]
    \centering
     \includegraphics[width=\textwidth]{figs/cloning_result.png}
     \caption{Seamless Cloning：边界处颜色平滑过渡，消除接缝，保持源图像细节纹理}
\end{figure}

\begin{figure}[H]
    \centering
     \includegraphics[width=\textwidth]{figs/mixing_result.png}
     \caption{Mixing Gradients：纹理融合更强，适用于保留背景纹理的场景}
\end{figure}

\subsubsection{Paste模式分析}
直接粘贴模式将源图像的像素值直接复制到目标图像中，存在以下问题：
\begin{itemize}
    \item 边界处颜色突变，产生明显的接缝。
    \item 光照条件不一致时，融合效果不自然。
    \item 无法保留目标图像的背景纹理。
\end{itemize}

\subsubsection{Seamless Cloning模式分析}
Seamless Cloning模式通过求解泊松方程实现无缝融合：
\begin{itemize}
    \item 边界处颜色平滑过渡，消除接缝。
    \item 保持源图像的细节纹理，同时适应目标图像的光照。
    \item 适用于克隆具有明显纹理的对象（如人脸、动物等）。
\end{itemize}

\subsubsection{Mixing Gradients模式分析}
Mixing Gradients模式在梯度域进行混合：
\begin{itemize}
    \item 在每个像素点选择梯度绝对值较大的源或目标梯度。
    \item 适用于克隆需要保留背景纹理的对象（如文字、水印等）。
    \item 在半透明或纹理丰富的背景下效果更佳。
\end{itemize}

\subsection{Mixing Gradients 效果分析（参数实验）}
\texttt{Mixing Gradients} 模式通过取源梯度和目标梯度中绝对值较大者作为引导场 $\mathbf{v}$。

\subsubsection{梯度选择策略}
对于每个像素点 $p$ 和每个方向 $d$，引导梯度定义为：
\begin{equation}
    v_{pd} = \begin{cases} \nabla f^*(p)_d, & \text{if } |\nabla f^*(p)_d| \geq |\nabla g(p)_d| \\ \nabla g(p)_d, & \text{otherwise} \end{cases}
\end{equation}
这种策略确保在目标图像梯度较强的区域保留背景纹理，在源图像梯度较强的区域保留源图像细节。

\subsubsection{实验场景分析}
\begin{itemize}
    \item \textbf{场景1：纯色背景}：Seamless Cloning和Mixing Gradients效果相近，因为目标梯度接近零。
    \item \textbf{场景2：纹理背景}：Mixing Gradients效果显著优于Seamless Cloning，能更好地保留背景纹理。
    \item \textbf{场景3：光照变化}：Seamless Cloning能更好地适应光照变化，Mixing Gradients可能产生不自然的梯度混合。
\end{itemize}

\subsubsection{定量分析}
通过计算融合区域与目标图像边界的平均梯度差，可以量化两种模式的差异：
\begin{itemize}
    \item Seamless Cloning：边界梯度差较小，过渡平滑。
    \item Mixing Gradients：边界梯度差较大，但保留了更多纹理信息。
\end{itemize}

\begin{figure}[H]
    \centering
    \includegraphics[width=1\textwidth]{figs/freeform_border.png}
    \caption{自由绘制边界实验：不规则选区下的 Mixing Gradients 结果，边缘过渡依然平滑。}
\end{figure}



\subsection{不规则区域实验}
\subsubsection{多边形选区实验}
使用多边形选区克隆复杂形状的对象：
\begin{itemize}
    \item 扫描线算法能准确识别多边形内部像素。
    \item 泊松方程求解能处理不规则边界。
    \item 边界过渡平滑，无明显接缝。
\end{itemize}

\subsubsection{自由绘制选区实验}
使用自由绘制选区克隆任意形状的对象：
\begin{itemize}
    \item 自由绘制产生大量顶点，扫描线算法仍能高效处理。
    \item 边界框优化减少了不必要的计算。
    \item 融合效果自然，边缘过渡平滑。
\end{itemize}

\subsubsection{边界处理策略}
对于不规则区域，采用动态邻居检测策略：
\begin{itemize}
    \item 检查四个方向是否在选中区域内。
    \item 在掩码外的邻居作为已知边界，移到方程右侧。
    \item 避免了强制截断为Dirichlet边界的问题。
\end{itemize}



\section{遇到的问题与解决方案}

\subsection{坐标映射错误}
\textbf{问题描述：} 出现“图像变白”与“掩码索引错误”的问题。
\begin{itemize}
    \item \textbf{原因：} 掩码像素索引计算 \texttt{(y * W + x) * channels} 时，未考虑目标图像相对于边界框的偏移。局部坐标（相对于边界框）和全局坐标（相对于源图像）混淆。
    \item \textbf{解决：} 修正了 \texttt{SeamlessClone::fill\_coefficient} 中的逻辑，通过引入 \texttt{bbox\_min} 确保局部坐标与全局像素位置的一一对应。在 \texttt{g()} 和 \texttt{f()} 函数中正确处理坐标转换。
\end{itemize}

\subsection{矩阵奇异问题}
\textbf{问题描述：} 求解器返回失败，提示矩阵可能奇异。
\begin{itemize}
    \item \textbf{原因：} 选中区域为空或区域形状导致矩阵秩亏。例如，单像素区域或线性区域。
    \item \textbf{解决：} 在 \texttt{precompute()} 中添加检查，如果选中像素数量为0，直接返回目标图像。同时，在UI中限制最小选区大小。
\end{itemize}

\subsection{边界处理问题}
\textbf{问题描述：} 不规则区域的边界处理不正确，导致融合结果异常。
\begin{itemize}
    \item \textbf{原因：} 最初将内部边缘点强制截断为Dirichlet边界，导致方程组结构不正确。
    \item \textbf{解决：} 采用动态邻居检测策略，对于掩码外的邻居作为已知边界，移到方程右侧。所有掩码内的点都使用Poisson方程。
\end{itemize}

\subsection{实时编辑卡顿}
\textbf{问题描述：} 开启实时编辑后，拖动时帧率较低，体验不佳。
\begin{itemize}
    \item \textbf{原因：} 每次拖动都重新分解矩阵，计算开销大。Debug模式下运行速度慢。
    \item \textbf{解决：} 实现矩阵预分解技术，只在首次点击时分解矩阵。后续拖动只更新右侧向量并执行回代。同时，确保在Release模式下运行。
\end{itemize}

\subsection{像素值溢出}
\textbf{问题描述：} 融合结果中某些像素值超出[0,255]范围，导致图像异常。
\begin{itemize}
    \item \textbf{原因：} 泊松方程的解可能为负数或超过255，直接作为像素值会导致截断。
    \item \textbf{解决：} 在 \texttt{solve()} 和 \texttt{solve\_fast()} 中，使用 \texttt{std::clamp()} 将解裁剪到[0,255]范围。
\end{itemize}

\subsection{混合梯度模式错误}
\textbf{问题描述：} Mixing Gradients模式下，融合结果不正确。
\begin{itemize}
    \item \textbf{原因：} 在计算混合梯度时，使用了旧的偏移量，导致目标图像梯度计算错误。
    \item \textbf{解决：} 在 \texttt{solve\_fast()} 中，每次拖动都使用最新的偏移量重新计算右侧向量B。确保混合梯度计算使用正确的目标图像位置。
\end{itemize}

\section{实验示例}

\subsection{示例1：人脸克隆}
\begin{figure}[H]
    \centering
    \includegraphics[width=0.45\textwidth]{figs/face_cloning.png}
    \caption{人脸克隆示例：将源图像中的人脸无缝克隆到目标图像中，边界过渡自然。}
\end{figure}
\begin{itemize}
    \item 源图像：包含清晰人脸在水中的照片。
    \item 目标图像：风景照片。
    \item 选区：矩形选区包围人脸。
    \item 结果：人脸无缝融合到风景中。
\end{itemize}

\subsection{示例2：文字水印}
\begin{figure}[H]
    \centering
    \includegraphics[width=0.45\textwidth]{figs/text_watermark.png}
    \caption{文字水印示例：使用Mixing Gradients模式将文字添加到纹理背景中。}
\end{figure}
\begin{itemize}
    \item 源图像：包含文字的透明背景图像。
    \item 目标图像：纹理丰富的背景图像。
    \item 选区：usct文字。
    \item 结果：文字自然地融入背景，同时保留了背景的纹理细节。
\end{itemize}


\section{结论}
本实验成功实现了基于泊松方程的图像无缝克隆。通过矩阵预分解技术，解决了大规模线性方程组求解带来的计算瓶颈，实现了流畅的实时交互。Mixing Gradients 模式的引入进一步拓宽了算法的应用场景，使其在处理纹理混合时表现出更强的鲁棒性。

\subsection{主要贡献}
\begin{enumerate}
    \item 实现了完整的Poisson Image Editing算法，包括Seamless Cloning和Mixing Gradients两种模式。
    \item 设计了高效的坐标映射机制，支持不规则区域的泊松方程求解。
    \item 实现了矩阵预分解技术，将实时编辑的帧率提升到60+ FPS。
    \item 扩展了区域选择功能，支持矩形、多边形和自由绘制三种选区类型。
\end{enumerate}

\subsection{实验总结}
\begin{itemize}
    \item 泊松方程能有效消除图像克隆时的边界接缝，实现无缝融合。
    \item 混合梯度模式在纹理丰富的背景下效果更佳，能更好地保留背景细节。
    \item 矩阵预分解技术是实现实时编辑的关键，性能提升显著。
    \item 不规则区域处理需要仔细设计坐标映射和边界条件。
\end{itemize}

\subsection{未来工作}
\begin{itemize}
    \item 实现更复杂的梯度场，如基于纹理的梯度引导。
    \item 优化大规模区域的求解性能，探索并行计算方法。
    \item 支持更多选区类型，如椭圆、贝塞尔曲线等。
    \item 实现交互式梯度编辑，允许用户手动调整引导场。
\end{itemize}

% 关闭中文环境
\end{CJK*}
\end{document}